\chapter{Analiza powierzchni ataku i scenariusze socjotechniczne}\label{ch:socjotechnika}

\section{Synteza powierzchni ataku i wektory podatności}\label{sec:synteza-powierzchni-ataku}
\subsection{Zestawienie kluczowych punktów ekspozycji D-INFK z danych OSINT}\label{subsec:punkty-ekspozycji}
\subsection{Podatności procedur organizacyjnych i czynnika ludzkiego}\label{subsec:podatnosci-ludzkie}

\section{Ramy operacyjne i rygorystyczne procedury bezpieczeństwa testów}\label{sec:ramy-bezpieczenstwa}
\subsection{Warunki bezpieczeństwa gwarantujące brak zakłóceń i szkód materialnych}\label{subsec:brak-szkod}
\subsection{Mechanizmy uniemożliwiające dostęp osobom postronnym (tokeny UUID, IP whitelisting)}\label{subsec:mechanizmy-ochronne}
\subsection{Zasada zerowej retencji haseł (Zero-Knowledge) i procedura natychmiastowego przerwania}\label{subsec:kill-switch}

\section{Scenariusz ataku nr 1: Kradzież poświadczeń z ominięciem 2FA}\label{sec:scenariusz-1}
\subsection{Grupa docelowa, pretexting i uzasadnienie w oparciu o OSINT}\label{subsec:scen1-cel}
\subsection{Szczegółowy opis techniczny (Evilginx2, domena typosquattingowa, makieta Switch edu-ID)}\label{subsec:scen1-techniczny}
\subsection{Zabezpieczenia techniczne przed dostępem osób postronnych}\label{subsec:scen1-zabezpieczenia}
\subsection{Analiza ryzyka operacyjnego i technicznego wraz z procedurami awaryjnymi}\label{subsec:scen1-ryzyko}

\section{Scenariusz ataku nr 2: Dostarczenie pliku z nieszkodliwym ładunkiem}\label{sec:scenariusz-2}
\subsection{Grupa docelowa, pretexting i uzasadnienie w oparciu o OSINT (pretekst grantu SNSF)}\label{subsec:scen2-cel}
\subsection{Szczegółowy opis techniczny (wektor PDF/DOCX, nieszkodliwy canary token)}\label{subsec:scen2-techniczny}
\subsection{Zabezpieczenia techniczne przed dostępem osób postronnych}\label{subsec:scen2-zabezpieczenia}
\subsection{Analiza ryzyka operacyjnego i technicznego wraz z procedurami awaryjnymi}\label{subsec:scen2-ryzyko}

\section{Scenariusz ataku nr 3: Podszycie się pod wsparcie techniczne helpdesku ISG}\label{sec:scenariusz-3}
\subsection{Grupa docelowa, pretexting i uzasadnienie w oparciu o OSINT}\label{subsec:scen3-cel}
\subsection{Szczegółowy opis techniczny (scenariusz rozmowy, wykorzystanie bazy wiedzy ISG, spoofing)}\label{subsec:scen3-techniczny}
\subsection{Zabezpieczenia techniczne przed dostępem osób postronnych}\label{subsec:scen3-zabezpieczenia}
\subsection{Analiza ryzyka operacyjnego i technicznego wraz z procedurami awaryjnymi}\label{subsec:scen3-ryzyko}

